% Generated draft: requires mathematical and source review.
\documentclass[11pt]{article}
\usepackage[margin=25mm]{geometry}
\usepackage{fontspec}
\setmainfont{Latin Modern Roman}
\usepackage{amsmath,amssymb,mathtools,unicode-math}
\setmathfont{Latin Modern Math}
\usepackage{enumitem,needspace,xurl,hyperref}
\setlength{\parindent}{0pt}
\setlength{\parskip}{5pt}
\setlength{\emergencystretch}{3em}
\newcommand{\N}{\mathbb{N}}
\newcommand{\Z}{\mathbb{Z}}
\newcommand{\Q}{\mathbb{Q}}
\newcommand{\R}{\mathbb{R}}
\newcommand{\C}{\mathbb{C}}
\newcommand{\F}{\mathbb{F}}
\newcommand{\A}{\mathbb{A}}
\newcommand{\PP}{\mathbb{P}}
\newcommand{\E}{\mathbb{E}}
\newcommand{\eps}{\varepsilon}
\newcommand{\dd}{\,\mathrm{d}}
\begin{document}
\section*{UnsolvedMath ID: 11000221}
\textbf{Source problem code:} catalogue entry
\subsection{Definitions and mathematical statement}
Fix a connected component \(\mathcal C\) of an abelian-differential stratum \(\mathcal H(\alpha)\), or of a quadratic-differential stratum \(\mathcal Q(\beta)\), with the source's zero/pole and marking conventions. A saddle connection is a flat geodesic segment joining singularities with no singularity in its interior. The problem is to determine the configuration types of homologous saddle connections, or homologous closed geodesics, occurring on \emph{every} surface in \(\mathcal C\).

If \(\mathscr T(X)\) is the set of admissible configuration types present on \(X\), the target is
\[
 \bigcap_{X\in\mathcal C}\mathscr T(X).
\]
Thus a type must be realized on each surface, although its direction and representative curves may vary; realization somewhere or almost everywhere is insufficient. Ordinary relative homology with the zero set is the usual abelian setting, after compatible orientations. For quadratic differentials the source expressly requires homology with local coefficients [MaZo], related to the sign local system for square roots of the quadratic differential; ordinary untwisted homology is not substituted. The exact extension at singularities, equivalence of configurations, maximality and primitive/multiplicity conventions are absent.

The continuation introduces the Teichm\"uller flow by the action of
\[
 \begin{pmatrix}e^t&0\\0&e^{-t}\end{pmatrix},\qquad t\in\R,
\]
on abelian strata. It points to Lyapunov exponents [BaPe, Fo2], reports Avila--Viana [AvVi] proving simplicity of the related cocycle spectrum, strengthening Forni's positivity result, and announces two later flow problems. The exact cocycle, invariant measure and convention for smallest positive versus smallest full-spectrum exponent are not given; no uniform theorem for arbitrary measures or missing subsequent problem is invented.
\subsection{Short English statement}
Find the homologous saddle-connection or closed-geodesic configuration types that occur on every surface of a fixed connected stratum component, using local-coefficient homology for quadratic differentials.
\subsection{Sources}
\begin{itemize}
\item \emph{Problems on Mapping Class Groups and Related Topics}, edited by Farb, Problem 14, ``Decomposition of surfaces,'' and following flow background, supplied PDF page 261. \url{https://www.math.uchicago.edu/~farb/papers/mcgbook.pdf}. Recorded access 2026-07-29; extraction pdf-font-structure.
\item [MaZo] for the local-coefficient homology convention; [BaPe], [Fo2] for Lyapunov background; A. Avila and M. Viana [AvVi] for reported simplicity, with earlier positivity attributed to Forni. Full bibliography entries and precise hypotheses are absent.
\item Local UnsolvedMath registry, \texttt{data/unsolvedmath/problems.json}, permanent ID \(11000221\), exact source code \texttt{catalogue entry}; supplied OPEN-TRIAGE summary reports unverified status.
\end{itemize}
\end{document}
